% Vietnamese translation for OI Public Library
% Source-File: IOI中国国家候选队论文/2004/韩文韬.ppt
\documentclass[11pt]{article}
\usepackage{oipl-vn}

\title{Ứng dụng C++ trong thi tin học\\{\large Ghi chú theo slide}}
\author{Han Wentao}
\date{IOI 2004}

\begin{document}
\maketitle

\origtitle{Using C++ in informatics contests}
\sourcepath{IOI 2004 / Han Wentao.ppt}

\section{Từ Pascal sang C++}

Slide mở bằng so sánh BASIC, Pascal, Java và C++. Trục chính của bài là giúp
người học Pascal chuyển sang C++: kiểu dữ liệu, tham chiếu, con trỏ, mảng,
cấu trúc, hàm và thư viện chuẩn. Trong thi đấu, lợi ích lớn của C++ không chỉ
là cú pháp, mà còn là khả năng trừu tượng hóa mà vẫn giữ hiệu năng tốt.

\section{Khuôn mẫu và lập trình tổng quát}

Phần trích được nhắc ví dụ \texttt{swap} với \texttt{template<class T>}, rồi
kiểu mảng cố định \texttt{c\_array<T,max>}. Đây là cánh cửa vào lập trình
tổng quát: viết một thuật toán cho cả họ kiểu dữ liệu, miễn các thao tác cần
thiết được định nghĩa.

\section{STL}

Các slide dành phần lớn dung lượng cho STL:
\begin{itemize}
  \item bộ lặp: input, output, forward, bidirectional và random access;
  \item \texttt{pair} và \texttt{make\_pair};
  \item thuật toán: \texttt{find}, \texttt{count}, \texttt{copy},
  \texttt{sort}, \texttt{lower\_bound}, \texttt{merge},
  \texttt{set\_union}, \texttt{make\_heap}, \texttt{min\_element};
  \item đối tượng hàm và predicate cho các thuật toán nhận tiêu chí so sánh.
\end{itemize}

\section{Ghi nhớ}

STL giúp giảm lỗi cài đặt các thao tác thường gặp, nhưng cần hiểu yêu cầu về
bộ lặp và độ phức tạp của từng thuật toán. Dùng đúng vật chứa và thuật toán
chuẩn là một kỹ năng thuật toán, không chỉ là mẹo ngôn ngữ.

\end{document}
